\section{Class NP --- Poly Time Prefix}
%%% generated by the blueprint pipeline from TCSlib.Complexity.ClassNP.PolyTimePrefix
\subsection*{Class NP --- Poly Time Prefix: Introduction}

Declarations of the Lean module \texttt{TCSlib.Complexity.ClassNP.PolyTimePrefix}.

\subsection*{Class NP --- Poly Time Prefix: Results}

\begin{definition}[Prefix bounded by decoded length]
\lean{Complexity.PrefixByLength.take}
\label{Complexity.PrefixByLength.take}
\leanok
\uses{Complexity.pairFstD, Complexity.pairSndD}
For a Boolean list $z$, decode it as a pair of lists and take from the second list the
prefix whose length is the length of the first list.
\end{definition}

\begin{definition}[Length-based list truncation]
\lean{Complexity.PrefixByLength.drop}
\label{Complexity.PrefixByLength.drop}
\leanok
\uses{Complexity.pairFstD, Complexity.pairSndD}
For a Boolean list $z$, form the list obtained from the second component of the
decomposition of $z$ and remove its first as many entries as the first component of that
decomposition has entries.
\end{definition}

\begin{definition}[Control labels for a prefix procedure]
\lean{Complexity.PrefixByLength.Label}
\label{Complexity.PrefixByLength.Label}
\leanok
Introduce a finite set of labels consisting of the states scan, second-false,
second-true, bump, head, read, decrement, copy, and stop, together with emit-$b$ and
copy-and-emit-$b$ for each Boolean bit $b$.
\end{definition}

\begin{definition}[Single-register initialization]
\lean{Complexity.PrefixByLength.regs}
\label{Complexity.PrefixByLength.regs}
\leanok
For each natural number $n$, assign the value $n$ to the sole register.
\end{definition}

\begin{definition}[Prefix-processing counter program]
\lean{Complexity.PrefixByLength.program}
\label{Complexity.PrefixByLength.program}
\leanok
\uses{Complexity.CounterProg.Instr, Complexity.PrefixByLength.Label}
For each choice of a Boolean parameter $k$, define a one-register counter program with
the following instructions at its labels. At scan, read a bit, stopping at end of input
and otherwise proceeding to second-false or second-true according to the bit. At
second-false, read a bit, stopping at end of input, proceeding to bump on $0$, and to
head on $1$; at second-true, read a bit, stopping at end of input, proceeding to stop on
$0$, and to bump on $1$. At bump, increment the register and return to scan. At head,
stop if the register is zero and $k$ is true, or proceed to copy if it is zero and $k$
is false; if the register is positive, proceed to read. At read, stop at end of input;
if $k$ is true, print the bit read and proceed to decrement, while if $k$ is false,
proceed to decrement for either bit. At decrement, decrement the register, leaving zero
unchanged, and return to head. At emit-$b$, print $b$ and proceed to decrement; at copy,
read a bit, stopping at end of input and otherwise proceeding to copy-and-emit that bit;
at copy-and-emit-$b$, print $b$ and return to copy; and at stop, halt.
\end{definition}

\begin{theorem}[Updating a Single Register]
\lean{Complexity.PrefixByLength.update_regs}
\label{Complexity.PrefixByLength.update_regs}
\leanok
\uses{Complexity.PrefixByLength.regs}
\difficulty{4}
For natural numbers $n$ and $m$, let the register state of size $n$ assign the value $n$
to its sole register. Updating that register to $m$ gives exactly the register state of
size $m$.
\end{theorem}

\begin{theorem}[Copying a Bit String]
\lean{Complexity.PrefixByLength.goes_copy}
\label{Complexity.PrefixByLength.goes_copy}
\leanok
\uses{Complexity.CounterProg.Goes, Complexity.CounterProg.Goes.congr, Complexity.CounterProg.Goes.prepend_input, Complexity.CounterProg.Goes.trans, Complexity.CounterProg.goes_halt, Complexity.CounterProg.goes_out, Complexity.CounterProg.goes_rd_end, Complexity.CounterProg.goes_rd_false, Complexity.CounterProg.goes_rd_true, Complexity.PrefixByLength.program, Complexity.PrefixByLength.regs}
\difficulty{5}
For either choice of the program parameter, every natural number $n$, and every bit
string $r$, on input $r$ the program starting at its copy label with its single register
equal to $n$ and input position $0$ reaches a halted state with that register unchanged,
input position $|r|$, and output extended by $r$, in at most $2|r|+2$ steps. This holds
for every output string already printed.
\end{theorem}

\begin{theorem}[Decrementing the Register Once]
\lean{Complexity.PrefixByLength.goes_dec_one}
\label{Complexity.PrefixByLength.goes_dec_one}
\leanok
\uses{Complexity.CounterProg.Goes, Complexity.CounterProg.goes_dec, Complexity.PrefixByLength.program, Complexity.PrefixByLength.regs, Complexity.PrefixByLength.update_regs}
\difficulty{3}
For either choice of the program parameter, every input bit string $x$, and all natural
numbers $n$ and $p$, the program starting at its decrement label with its single
register equal to $n+1$ and input position $p$ reaches its head label after one step,
with register value $n$ and input position still $p$, adding no output. This holds for
every output string already printed.
\end{theorem}

\begin{theorem}[Processing the Payload Prefix or Suffix]
\lean{Complexity.PrefixByLength.goes_payload}
\label{Complexity.PrefixByLength.goes_payload}
\leanok
\uses{Complexity.CounterProg.Goes, Complexity.CounterProg.Goes.congr, Complexity.CounterProg.Goes.prepend_input, Complexity.CounterProg.Goes.trans, Complexity.CounterProg.goes_halt, Complexity.CounterProg.goes_jz_pos, Complexity.CounterProg.goes_jz_zero, Complexity.CounterProg.goes_out, Complexity.CounterProg.goes_rd_end, Complexity.CounterProg.goes_rd_false, Complexity.CounterProg.goes_rd_true, Complexity.PrefixByLength.goes_copy, Complexity.PrefixByLength.goes_dec_one, Complexity.PrefixByLength.program, Complexity.PrefixByLength.regs}
\difficulty{6}
For either choice of the program parameter, every bit string $r$, and every natural
number $n$, the program on input $r$, starting at its head label with register value $n$
and input position $0$, reaches a halted state with register value $n-|r|$ truncated at
zero. Its final input position is $\min(n,|r|)$ when the parameter is true and $|r|$
when it is false; it appends the first $n$ bits of $r$ when the parameter is true and
the suffix of $r$ after its first $n$ bits when false. The run takes at most $4|r|+3$
steps, for every output string already printed.
\end{theorem}

\begin{theorem}[Incrementing the Register Once]
\lean{Complexity.PrefixByLength.goes_inc_one}
\label{Complexity.PrefixByLength.goes_inc_one}
\leanok
\uses{Complexity.CounterProg.Goes, Complexity.CounterProg.goes_inc, Complexity.PrefixByLength.program, Complexity.PrefixByLength.regs, Complexity.PrefixByLength.update_regs}
\difficulty{3}
For either choice of the program parameter, every input bit string $x$, and all natural
numbers $n$ and $p$, the program starting at its bump label with its single register
equal to $n$ and input position $p$ reaches its scan label after one step, with register
value $n+1$ and input position still $p$, adding no output. This holds for every output
string already printed.
\end{theorem}

\begin{theorem}[Scanning a Doubled Prefix]
\lean{Complexity.PrefixByLength.goes_doubled}
\label{Complexity.PrefixByLength.goes_doubled}
\leanok
\uses{Complexity.CounterProg.Goes, Complexity.CounterProg.Goes.congr, Complexity.CounterProg.Goes.prepend_input, Complexity.CounterProg.Goes.trans, Complexity.CounterProg.goes_rd_false, Complexity.CounterProg.goes_rd_true, Complexity.CounterProg.run_zero, Complexity.PrefixByLength.goes_inc_one, Complexity.PrefixByLength.program, Complexity.PrefixByLength.regs, Turing.dbl, Turing.dbl_nil}
\difficulty{4}
For either choice of the program parameter, all bit strings $u$ and $t$, and every
natural number $n$, on input formed by writing each bit of $u$ twice and then appending
$t$, the program starting at its scan label with register value $n$ and input position
$0$ reaches the scan label with register value $n+|u|$ and input position $2|u|$, adding
no output, in at most $3|u|$ steps. This holds for every output string already printed.
\end{theorem}

\begin{theorem}[Processing a Self-Delimited Pair]
\lean{Complexity.PrefixByLength.goes_pair}
\label{Complexity.PrefixByLength.goes_pair}
\leanok
\uses{Complexity.CounterProg.Goes, Complexity.CounterProg.Goes.congr, Complexity.CounterProg.Goes.prepend_input, Complexity.CounterProg.Goes.trans, Complexity.CounterProg.goes_rd_false, Complexity.CounterProg.goes_rd_true, Complexity.PrefixByLength.goes_doubled, Complexity.PrefixByLength.goes_payload, Complexity.PrefixByLength.program, Complexity.PrefixByLength.regs, Turing.dbl, Turing.length_pairEncode, Turing.pairEncode, Turing.pairEncode_eq_dbl}
\difficulty{5}
For either choice of the program parameter and all bit strings $u$ and $r$, there are a
register assignment and an input position such that, on the input formed by doubling
every bit of $u$, appending the separator $01$, and then appending $r$, the program
starting at its scan label with register value $0$ and input position $0$ reaches a
halted state with that register assignment and input position. It adds $r$ truncated to
its first $|u|$ bits when the parameter is true, and $r$ with its first $|u|$ bits
removed when the parameter is false; the run takes at most five times one more than the
input length. This holds for every output string already printed.
\end{theorem}

\begin{theorem}[Halting on a Short or Delimited Tail]
\lean{Complexity.PrefixByLength.goes_bad_tail}
\label{Complexity.PrefixByLength.goes_bad_tail}
\leanok
\uses{Complexity.CounterProg.Goes, Complexity.CounterProg.Goes.congr, Complexity.CounterProg.Goes.trans, Complexity.CounterProg.goes_halt, Complexity.CounterProg.goes_rd_end, Complexity.CounterProg.goes_rd_false, Complexity.CounterProg.goes_rd_true, Complexity.PrefixByLength.program, Complexity.PrefixByLength.regs}
\difficulty{4}
For either choice of the program parameter and every natural number $n$, if a bit string
$t$ is empty, has length one, or begins with $10$, then there is a natural number $p$
such that, on input $t$, the program starting at its scan label with its single register
equal to $n$ and input position $0$ reaches a halted state with the same register value
and input position $p$, without adding any output, in at most three steps. This holds
for every output string already printed.
\end{theorem}

\begin{theorem}[Total Processing of a Pair Encoding]
\lean{Complexity.PrefixByLength.goes_total}
\label{Complexity.PrefixByLength.goes_total}
\leanok
\uses{Complexity.CounterProg.Goes, Complexity.CounterProg.Goes.congr, Complexity.CounterProg.Goes.prepend_input, Complexity.CounterProg.Goes.trans, Complexity.PrefixByLength.drop, Complexity.PrefixByLength.goes_bad_tail, Complexity.PrefixByLength.goes_doubled, Complexity.PrefixByLength.goes_pair, Complexity.PrefixByLength.program, Complexity.PrefixByLength.regs, Complexity.PrefixByLength.take, Complexity.pairFstD, Complexity.pairSndD, Turing.dbl, Turing.eq_pairEncode_of_pairDecode, Turing.pairDecode, Turing.pairDecode_eq_none}
\difficulty{5}
For either choice of the program parameter and every bit string $z$, there are a
register assignment and an input position such that, on input $z$, the program starting
at its scan label with register value $0$ and input position $0$ reaches a halted state
with that assignment and position. To describe its appended output, parse $z$ into a
first component $a$ and a second component $b$ by reading aligned pairs $00$ and $11$ as
bits of $a$ until the separator $01$, with the remaining suffix as $b$; if this parse
fails, take both components to be empty. The output appended is the first $|a|$ bits of
$b$ when the parameter is true and the suffix of $b$ after its first $|a|$ bits when
false, and the run takes at most $5(|z|+1)$ steps; this holds for every output string
already printed.
\end{theorem}

\begin{theorem}[Polynomial-time length-bounded prefix]
\lean{Complexity.polyTimeComputable_takePrefixByLength}
\label{Complexity.polyTimeComputable_takePrefixByLength}
\leanok
\uses{Complexity.CounterProg.polyTimeComputable_of_goes, Complexity.PolyTimeComputable, Complexity.PrefixByLength.goes_total, Complexity.PrefixByLength.program, Complexity.PrefixByLength.take}
\difficulty{1}
For a bit string $z$, parse its two components using the aligned pair parser, taking
both components to be empty if parsing fails. The function that returns the prefix of
the second component whose length is the length of the first component is
polynomial-time computable.
\end{theorem}

\begin{theorem}[Dropping a length-sized prefix is polynomial-time]
\lean{Complexity.polyTimeComputable_dropPrefixByLength}
\label{Complexity.polyTimeComputable_dropPrefixByLength}
\leanok
\uses{Complexity.CounterProg.polyTimeComputable_of_goes, Complexity.PolyTimeComputable, Complexity.PrefixByLength.drop, Complexity.PrefixByLength.goes_total, Complexity.PrefixByLength.program}
\difficulty{1}
For a bit string $z$, let $a$ and $b$ be the two components decoded from $z$ by the
aligned pair parser, using the empty word for each total projection when decoding fails.
The function that drops from $b$ its first $|a|$ bits is polynomial-time computable.
\end{theorem}
